\section{结果分析}\label{sec:task2-analysis}

本节采用“实测数值---判据---结论”的顺序填写。由于真实照片和标定输出尚未在本报告中落地，以下内容是待数据补入后直接执行的分析框架，不预先声称标定成功或误差达到某一阈值。

\subsection{内参合理性}

首先检查 \(f_x\)、\(f_y\)、\(c_x\) 和 \(c_y\) 是否为有限的正数，并将其与图像尺寸和式\ref{eq:task2-intrinsic}逐项核对。记录
\[
\rho_f=\frac{\lvert f_x-f_y\rvert}{(f_x+f_y)/2}
\]
作为两个方向焦距差异的相对量；\(\rho_f\) 的具体数值及解释应结合手机镜头、图像缩放和采集设置，不应脱离数据直接规定结论。主点与图像几何中心的偏移可由
\[
\Delta_c=\sqrt{(c_x-W/2)^2+(c_y-H/2)^2}
\]
计算，其中 \(W,H\) 是实际图像宽、高。应报告 \(\Delta_c\) 及其相对于图像尺寸的比例，而不是仅凭“接近中心”作定性判断。

\subsection{畸变与误差}

按 OpenCV 的实际输出顺序记录 \(k_1,k_2,p_1,p_2,k_3,\ldots\)，结合原始图像边缘和校正图像观察径向变化及切向偏差。系数正负本身不能单独证明校正效果，需与校正前后直线或棋盘格边缘的变化对应。将 OpenCV RMS 与平均重投影误差分别和式\ref{eq:task2-point-reprojection}、式\ref{eq:task2-reprojection-error}核对；若二者差异异常，应先检查误差统计口径、角点数组形状和单位。

\subsection{数据质量与姿态覆盖}

根据表\ref{tab:task2-data-record}的采集记录核对标定板是否覆盖中央、边缘和不同倾角。结合每幅图像的角点检测状态，报告有效图像数、失败图像数、失败原因及单幅误差分布。若某些照片模糊、反光、遮挡或姿态高度重复，应说明其可能造成的影响，并在需要时剔除后重新运行，比较两次结果而不是直接删除异常值。

\begin{table}[htbp]
    \centering
    \caption{结果分析记录框架}
    \label{tab:task2-analysis}
    \begin{tabular}{lll}
        \hline
        检查项 & 需填入的证据 & 分析结论 \\
        \hline
        焦距一致性 & \(\rho_f=\) 待实际运行填入 & 待结合数据填写 \\
        主点位置 & \(\Delta_c=\) 待实际运行填入 & 待结合图像尺寸填写 \\
        畸变系数 & 实际完整向量 & 待结合校正图像填写 \\
        RMS 与平均误差 & 两个实际误差值 & 待结合角点质量填写 \\
        姿态覆盖 & 位置、倾角和距离记录 & 待根据采集记录填写 \\
        失败样本 & 文件名与失败原因 & 待根据程序日志填写 \\
        \hline
    \end{tabular}
\end{table}

只有当参数文件、程序日志、原始照片和表格中的配置值能够相互对应时，才可在本节给出“标定质量满足实验要求”等结论。若有效图像数量不足、图像尺寸混用或角点定位明显不稳定，应先补充数据或修正采集流程。
